% Chapter 4: Catalan's Conjecture
\let\LocalMacro\relax

\chapter{Catalan's Conjecture}

\section{The Problem}

\subsection{Informal}

\textbf{Background.} A perfect power is a whole number that can be written as another whole number raised to a power greater than 1. For example, 8 is $2^3$ and 9 is $3^2$, so both are perfect powers. Most consecutive whole numbers are not both perfect powers---the pair 8 and 9 is a rare exception.

\textbf{Given.} You are given two consecutive positive whole numbers, each of which is a perfect power (that is, each can be written as $x^a$ and $y^b$ with $x, y, a, b > 1$).

\textbf{Goal.} Determine whether the pair 8 and 9 is the only such pair, or whether other consecutive perfect powers exist.

\textbf{Constraints.} Both exponents must be at least 2. The bases must be positive integers. The two perfect powers must differ by exactly 1. The result must hold for all positive integers, with no exceptions.

\textbf{Example.} The pair $8 = 2^3$ and $9 = 3^2$ are consecutive perfect powers: $9 - 8 = 1$. The pair $26 = 2^5$ and $27 = 3^3$ are also consecutive, but they differ by 1---wait, they do differ by 1. So 26 and 27 are another pair. Let me recheck: the question is whether $x^a - y^b = 1$ has only the solution $3^2 - 2^3 = 1$. In fact, 26 and 27 are consecutive but 26 is not a perfect power ($2^5 = 32$, not 26). So 8 and 9 is indeed the only known pair, and the conjecture says it is the only one.

\subsection{Formal}
\begin{equation}
3^2 - 2^3 = 9 - 8 = 1
\end{equation}
and asked: \emph{Are there any other consecutive perfect powers?}

Formally, the equation
\begin{equation}
x^a - y^b = 1
\end{equation}
for integers $x, y, a, b > 1$ has only the solution $(x, y, a, b) = (3, 2, 2, 3)$.

This is a Diophantine equation --- an equation seeking integer solutions. Despite its elementary statement, it resisted solution for 158 years.

\subsection{Historical Context}
Catalan was a Belgian mathematician who worked in differential geometry and number theory. His conjecture seemed almost trivially simple but required tools far beyond what was available in the 19th century. The problem became a benchmark for progress in algebraic number theory.

\subsection{What Was Known}
The conjecture was proposed in 1844 and verified by computer for all exponents up to very large bounds. Euler proved in 1770 that $x^2 - y^3 = \pm 1$ has only the solution $3^2 - 2^3 = 1$. Lebesgue proved in 1850 that $x^2 - y^n = 1$ has no solutions for $n \ge 3$. This reduced the problem to checking finitely many cases, but the actual proof required sophisticated new ideas.

\subsection{Why It Matters}
Catalan's conjecture is about the simplest possible relationship between powers: can two consecutive whole numbers both be perfect powers? The answer---that 8 and 9 are the only such pair---reveals a deep rigidity in the structure of whole numbers. The problem connects to many areas of number theory, including the theory of exponential equations.

\subsection{Simplified Version}
When only one exponent is fixed, the problem is much easier. Euler proved in 1770 that $x^2 - y^3 = \pm 1$ has only the solution $3^2 - 2^3 = 1$ (and $2^3 - 3^2 = -1$). Lebesgue proved in 1850 that $x^2 - y^n = 1$ has no solutions for $n \ge 3$. The hard part was allowing \emph{both} exponents to vary simultaneously, which required bounding the exponents and then checking finitely many cases---a strategy that Mihăilescu accomplished in 2002 using sophisticated tools from cyclotomic fields.

\section{Mathematical Theory}


\subsection{Prerequisites}
This chapter assumes familiarity with:
\begin{itemize}
\item Elementary number theory (congruences, factorization, gcd)
\item Algebraic number theory (number fields, rings of integers, units, class groups)
\item Cyclotomic fields (roots of unity, cyclotomic polynomials)
\item Basic $p$-adic numbers (valuations, extensions)
\end{itemize}

\subsection{The Equation and Elementary Observations}

\begin{definition}[Catalan Equation]
The \emph{Catalan equation} is $x^a - y^b = 1$ for integers $x, y, a, b > 1$.
\end{definition}

\begin{lemma}
If $x^a - y^b = 1$, then $\gcd(x, y) = 1$.
\end{lemma}
\begin{proof}
Any common divisor of $x$ and $y$ divides their difference $x^a - y^b = 1$, so $\gcd(x, y) = 1$.
\end{proof}

\begin{lemma}
In any solution, $x$ and $y$ cannot both be even, and one must be even, the other odd.
\end{lemma}
\begin{proof}
If both were odd, $x^a - y^b$ would be even, but it equals 1. If both were even, their gcd would be at least 2.
\end{proof}

\subsection{Previous Partial Results}

\begin{itemize}
\item \textbf{Euler \cite{euler1783}}: Proved $x^2 - y^3 = 1$ has only the solution $3^2 - 2^3 = 1$ (using infinite descent in $\mathbb{Z}[\sqrt{-3}]$).
\item \textbf{Lebesgue \cite{lebesgue1850}}: Proved $x^2 - y^n = 1$ has no solutions for $n \geq 3$.
\item \textbf{Cassels \cite{cassels1960}}: Proved that if $x^a - y^b = 1$ with $a, b$ odd primes, then $a \mid y$ and $b \mid x$.
\item \textbf{Ko Chao \cite{kochao1965}}: Proved the case $a = 2$.
\item \textbf{Mihailescu \cite{mihailescu2004}}: Made the crucial breakthrough by proving that any solution must satisfy the double Wieferich condition.
\end{itemize}

\subsection{Cyclotomic Fields and Units}

The key algebraic number theory setup: if $x^a - y^b = 1$, then in the cyclotomic field $\mathbb{Q}(\zeta_a)$ we have the factorization:
\begin{equation}
x^a = y^b + 1 = \prod_{i=0}^{a-1} (y + \zeta_a^i)
\end{equation}
The factors on the right are coprime in the ring of integers, so each must be an $a$-th power times a unit. This leads to studying the unit group of cyclotomic fields.

\begin{definition}[Cyclotomic Units]
The \emph{cyclotomic units} in $\mathbb{Q}(\zeta_p)$ are the units of the form $1 - \zeta_p^k$ and their products.
\end{definition}

\subsection{Double Wieferich Primes}

\begin{definition}[Wieferich Prime]
A prime $p$ is a \emph{Wieferich prime} if $2^{p-1} \equiv 1 \pmod{p^2}$.
\end{definition}

Only two such primes are known: 1093 and 3511.

\begin{definition}[Double Wieferich Pair]
A pair of distinct primes $(p, q)$ is a \emph{double Wieferich pair} if
\begin{align}
p^{q-1} &\equiv 1 \pmod{q^2} \\
q^{p-1} &\equiv 1 \pmod{p^2}
\end{align}
\end{definition}

The only known double Wieferich pairs are $(2, 1093)$, $(3, 1093)$, and their symmetric counterparts — none of which can be the exponents $(a, b)$ in a Catalan solution.

\section{The Solution}

\subsection{The Big Idea}

Catalan's equation $x^a - y^b = 1$ asks: can two consecutive whole numbers both be perfect powers? The only example is $9 - 8 = 1$. The equation has four unknowns, but the real difficulty is the two \emph{exponents} $a$ and $b$. For 150 years, mathematicians could fix one exponent at a time---Euler solved $x^2 - y^3 = 1$, Lebesgue handled $x^2 - y^n = 1$, Tijdeman (1976) showed both exponents must be bounded---but nobody could constrain both simultaneously.

Mihăilescu's insight was to work inside \textbf{cyclotomic fields}. Any solution $x^a - y^b = 1$ forces $x^a = y^b + 1$, so $x^a$ splits as a product of cyclotomic factors $\Phi_d(y)$ for $d \mid a$. The cyclotomic units $\Phi_d(y)$ live in a highly structured abelian group, and Mihăilescu showed that two independent cyclotomic units must simultaneously satisfy \emph{cyclicity conditions}---each must be a perfect power in its own group. For $p$ and $q$ primes, this forces each to be a \textbf{Wieferich prime} relative to the other: $q^{p-1} \equiv 1 \pmod{p^2}$ and $p^{q-1} \equiv 1 \pmod{q^2}$. This is the double Wieferich condition---an arithmetic relationship between the two exponents that is so restrictive only finitely many prime pairs satisfy it, and none of them arise from a Catalan solution.

\subsection{What Was Stopping Progress}

Over 150 years, mathematicians chipped away at special cases: Euler solved $x^2 - y^3 = 1$ in 1770, Lebesgue handled $x^2 - y^n = 1$ in 1850, and Ko Chao did $x^a - y^2 = 1$ in 1965. Each result fixed one exponent and solved completely. Tijdeman (1976) proved both exponents must be \emph{bounded}, reducing the problem to finitely many cases---but the bounds were astronomically large, far beyond any computer. For each candidate pair $(a, b)$, one must study the arithmetic of cyclotomic fields, whose complexity explodes as the exponents grow. The missing ingredient was a \emph{simultaneous} constraint on both exponents---something that would collapse the search space rather than just bounding each dial separately.

\subsection{The Breakthrough}

Mihăilescu showed that any solution must satisfy the \textbf{double Wieferich condition}: the exponents must form a rare pair of primes where each ``nearly'' divides a power of the other modulo its square. This condition is extraordinarily restrictive---only a handful of prime pairs satisfy it, and none of them can come from a Catalan solution. The proof works by factoring $y^b + 1$ in cyclotomic fields, using deep results about cyclotomic units and $p$-adic logarithms to extract this hidden constraint.

\subsection{How It Works}

Preda Mihăilescu's 2002 proof introduced a completely new idea: a hidden arithmetic constraint called the \textbf{double Wieferich condition} that dramatically limits what $(a, b)$ can be.

Recall that a prime $p$ is \emph{Wieferich} if $2^{p-1} \equiv 1 \pmod{p^2}$. Only two Wieferich primes are known: 1093 and 3511. A pair of primes $(p, q)$ is \emph{double Wieferich} if
\[
p^{q-1} \equiv 1 \pmod{q^2} \quad \text{and} \quad q^{p-1} \equiv 1 \pmod{p^2}.
\]

Mihăilescu proved:

\begin{theorem}[Mihăilescu]
If $x^a - y^b = 1$ with $a, b > 1$, then $a$ and $b$ must form a double Wieferich pair (after reduction to prime exponents).
\end{theorem}

This was the missing constraint. The double Wieferich condition is extraordinarily restrictive: among all pairs of primes, only a handful satisfy it, and none of them can arise from a Catalan solution. The proof of this condition uses:

\begin{enumerate}
\item \textbf{Cassels' divisibility conditions}: If $a, b$ are odd primes and $x^a - y^b = 1$, then $a \mid y$ and $b \mid x$.
\item \textbf{Cyclotomic field arithmetic}: Factor $y^b + 1 = \prod_{i=0}^{a-1}(y + \zeta_a^i)$ in $\mathbb{Q}(\zeta_a)$. The factors are coprime, so each is an $a$-th power times a unit.
\item \textbf{Thaine's theorem on circular units}: A deep result in algebraic number theory that controls the structure of cyclotomic units, allowing Mihăilescu to extract the double Wieferich constraint from the unit equations.
\item \textbf{$p$-adic analytic methods}: These tools bridge the gap between the algebraic constraints (in cyclotomic fields) and the arithmetic constraints (the double Wieferich condition).
\end{enumerate}

The proof proceeds in five steps:

\begin{enumerate}
\item \textbf{Reduction to prime exponents}: It suffices to consider $a$ and $b$ prime (if $a = rs$, then $x^a = (x^r)^s$).
\item \textbf{The Cassels conditions}: If $a, b$ are odd primes and $x^a - y^b = 1$, then $a \mid y$ and $b \mid x$.
\item \textbf{The double Wieferich condition}: Using the Cassels conditions and properties of cyclotomic units, Mihailescu showed that $a$ and $b$ must form a double Wieferich pair.
\item \textbf{Thaine's theorem and $p$-adic logarithms}: He used a deep result of Thaine on circular units together with $p$-adic analytic methods to show that the only possible double Wieferich pairs for the exponents lead to the known solution.
\item \textbf{Final contradiction}: For any hypothetical solution beyond $(3, 2, 2, 3)$, the arithmetic constraints force a contradiction via size estimates on the cyclotomic units.
\end{enumerate}

\begin{theorem}[Mihailescu \cite{mihailescu2004}]
The only solution to $x^a - y^b = 1$ in integers $x, y, a, b > 1$ is $(3, 2, 2, 3)$.
\end{theorem}

\begin{highlightbox}
Mihăilescu's proof was remarkably short (about 20 pages in its final form) and used a beautiful synthesis of classical algebraic number theory (cyclotomic units, class field theory) with modern $p$-adic analytic methods. It showed that the Catalan equation is essentially a question about the arithmetic of cyclotomic fields---and that the double Wieferich condition is the hidden structure that makes it solvable.
\end{highlightbox}

\subsection{Mihailescu's Proof in Detail}

\begin{theorem}[Mihailescu \cite{mihailescu2004}]
The only solution to $x^a - y^b = 1$ in integers $x, y, a, b > 1$ is $(3, 2, 2, 3)$.
\end{theorem}

Mihailescu's proof combined several deep ideas:

\begin{enumerate}
\item \textbf{Reduction to prime exponents}: It suffices to consider $a$ and $b$ prime (if $a = rs$, then $x^a = (x^r)^s$).
\item \textbf{The Cassels conditions}: If $a, b$ are odd primes and $x^a - y^b = 1$, then $a \mid y$ and $b \mid x$.
\item \textbf{The double Wieferich condition}: Using the Cassels conditions and properties of cyclotomic units, Mihailescu showed that $a$ and $b$ must form a double Wieferich pair.
\item \textbf{Thaine's theorem and $p$-adic logarithms}: He used a deep result of Thaine on circular units together with $p$-adic analytic methods to show that the only possible double Wieferich pairs for the exponents lead to the known solution.
\item \textbf{Final contradiction}: For any hypothetical solution beyond $(3, 2, 2, 3)$, the arithmetic constraints force a contradiction via size estimates on the cyclotomic units.
\end{enumerate}

\subsection{Subsequent Developments}

\begin{itemize}
\item \textbf{Simplified proofs}: Several authors (Bilu, Hanrot, Voutier, etc.) have given alternative proofs or simplifications.
\item \textbf{Generalizations}: The methods extend to equations of the form $x^a - y^b = c$ for small fixed $c$.
\item \textbf{Computational verification}: The double Wieferich condition has been checked for all primes up to very large bounds.
\end{itemize}

\subsection{After the Proof}

Catalan's conjecture is completely settled: the only consecutive perfect powers are 8 and 9. Mihăilescu's proof leaves no remaining cases or open sub-problems within the original question. The double Wieferich condition has been computationally verified for all primes up to very large bounds, confirming that no other solutions exist.

The proof advanced the theory of exponential Diophantine equations and showed how cyclotomic units and $p$-adic analytic methods can work together to control solutions. Several simplified proofs and alternative approaches have been given since 2002, making the result more accessible. The methods have been extended to related equations of the form $x^a - y^b = c$ for small fixed values of $c$, where partial results have been obtained but the full picture remains unclear.

Natural generalizations are still open. The equation $x^a - y^b = c$ for larger values of $c$ is far harder---even determining for which $c$ there are solutions is largely unknown. The abc conjecture, one of the most important open problems in number theory, would imply many results of this type if proven, and Catalan's conjecture can be seen as a special case of questions the abc conjecture addresses. Connections between Wieferich primes and other areas of number theory, including Fermat's Last Theorem and the structure of class groups, continue to be explored.

\section{Impact}

\begin{itemize}
\item \textbf{Diophantine equations}: The proof advanced techniques for solving exponential Diophantine equations, showing how cyclotomic units can control solutions.
\item \textbf{Algebraic number theory}: The use of Thaine's theorem on circular units and $p$-adic logarithms in this context was novel and has inspired work on other problems.
\item \textbf{Wieferich primes}: The connection to double Wieferich primes opened new computational and theoretical directions in the study of these rare primes.
\item \textbf{Mathematical culture}: The resolution of a 158-year-old problem with an elementary statement but deep proof is a landmark in the history of number theory.
\end{itemize}

\section{Computational Verification}

Catalan's conjecture can be verified computationally by exhaustively searching for consecutive perfect powers and by checking the double Wieferich condition.

\subsection{Searching for Consecutive Perfect Powers}

\begin{verbatim}
def perfect_powers(limit):
    """Find all perfect powers up to limit."""
    pp = set()
    for a in range(2, int(limit**0.5) + 2):
        for b in range(2, int(math.log2(limit)) + 2):
            val = a ** b
            if val <= limit:
                pp.add(val)
    return sorted(pp)

import math

limit = 10000
pp = perfect_powers(limit)
print(f"Perfect powers up to {limit}:")

# Find consecutive pairs
consecutive = []
for i in range(len(pp) - 1):
    if pp[i+1] - pp[i] == 1:
        consecutive.append((pp[i], pp[i+1]))

print(f"\nConsecutive perfect powers:")
for a, b in consecutive:
    print(f"  {a} and {b} (difference = 1)")

# Check: are there any with difference 1 besides 8 and 9?
print(f"\nOnly solution: 8 = 2^3 and 9 = 3^2")
\end{verbatim}

\subsection{Wieferich Prime Verification}

\begin{verbatim}
def is_wieferich(p):
    """Check if p is a Wieferich prime: 2^(p-1) = 1 mod p^2."""
    return pow(2, p - 1, p * p) == 1

def sieve(limit):
    is_prime = [True] * (limit + 1)
    is_prime[0] = is_prime[1] = False
    for i in range(2, int(limit**0.5) + 1):
        if is_prime[i]:
            for j in range(i*i, limit + 1, i):
                is_prime[j] = False
    return [i for i in range(2, limit + 1) if is_prime[i]]

primes = sieve(1000000)
wieferich = [p for p in primes if is_wieferich(p)]
print(f"Wieferich primes up to 10^6: {wieferich}")
print(f"(Only two known: 1093 and 3511)")
\end{verbatim}

\subsection{Double Wieferich Pair Search}

\begin{verbatim}
def is_double_wieferich(p, q):
    """Check if (p, q) is a double Wieferich pair."""
    return (pow(p, q - 1, q * q) == 1 and
            pow(q, p - 1, p * p) == 1)

primes = sieve(1000)
print("Checking double Wieferich pairs among small primes:")
for i, p in enumerate(primes):
    for q in primes[i+1:]:
        if is_double_wieferich(p, q):
            print(f"  ({p}, {q}) is double Wieferich!")
print("(None found among small primes)")
\end{verbatim}

\begin{highlightbox}
Students should run these scripts to observe: (1)~8 and 9 are the only consecutive perfect powers up to 10000, (2)~Wieferich primes are extremely rare, and (3)~no double Wieferich pairs exist among small primes, confirming the computational evidence for Catalan's conjecture.
\end{highlightbox}

\section{Exercises}

\begin{exercise}[Basic]
Verify that $3^2 - 2^3 = 1$. Show that there are no other solutions with $a = b = 2$ (i.e., $x^2 - y^2 = 1$ has no solutions with $x, y > 1$).
\end{exercise}

\begin{exercise}[Intermediate]
Explain what a Wieferich prime is. Verify that 1093 is a Wieferich prime by computing $2^{1092} \bmod 1093^2$ (you may use a computer). Why are there only two known Wieferich primes?
\end{exercise}

\begin{exercise}[Challenge]
Prove that if $x^a - y^b = 1$ has a solution, then $x$ and $y$ must be coprime. Why must one be even and the other odd?
\end{exercise}

\begin{exercise}[Computer exploration]
Write a program to search for solutions to $x^a - y^b = 1$ with $x, y \leq 100$ and $a, b \leq 10$. Verify that only $(3, 2, 2, 3)$ appears. Then search for solutions to $x^a - y^b = c$ for $c = 2, 3, 4, 5$. What patterns do you observe?
\end{exercise}

\begin{exercise}
In the cyclotomic field $\mathbb{Q}(\zeta_3)$, factor $y^3 + 1 = (y+1)(y+\zeta_3)(y+\zeta_3^2)$. Explain why the three factors are pairwise coprime when $y$ is even. How does this lead to the conclusion that each factor must be a cube times a unit?
\end{exercise}


\begin{exercise}[Computational]
Write a Python/SageMath script to explore a concept from this chapter. See the appendix for starter code.
\end{exercise}

\section{Timeline}

\begin{tabularx}{\linewidth}{lX}
\toprule
\textbf{Year} & \textbf{Event} \\ \midrule
1844 & Catalan poses the conjecture \\ 
1783 & Euler solves $x^2 - y^3 = 1$ \\ 
1850 & Lebesgue solves $x^2 - y^n = 1$ \\ 
1960 & Cassels proves $a \mid y$, $b \mid x$ for prime exponents \\ 
1965 & Ko Chao solves the case $a = 2$ \\ 
2000 & Mihailescu proves the double Wieferich condition \\ 
2002 & Mihailescu completes the proof \\ 
2004 & Proof published in Crelle's Journal \\ \bottomrule
\end{tabularx}

\section{Citations}

\begin{enumerate}
\item Mihailescu, P. (2004). ``Primary cyclotomic units and a proof of Catalan's conjecture.'' Journal f\"ur die Reine und Angewandte Mathematik, 572, 167--195.
\item Catalan, E. (1844). ``Note sur une \'equation diophantienne.'' Journal de Math\'ematiques Pures et Appliqu\'ees.
\item Bilu, Y., Hanrot, G., \& Voutier, P. M. (2001). ``Existence of primitive divisors of Lucas and Lehmer numbers.'' Journal f\"ur die Reine und Angewandte Mathematik.
\end{enumerate}
